\documentclass{amsart}
\usepackage{amsmath,amssymb,amsthm,hyperref}

\theoremstyle{definition}
\newtheorem{definition}{Definition}
\newtheorem{thesis}{Thesis}
\theoremstyle{plain}
\newtheorem{claim}{Claim}
\newtheorem{lemma}{Lemma}
\newtheorem{proposition}{Proposition}
\newtheorem{objection}{Objection}

\begin{document}

\title{Scientific Objectivity and the Separation of the Descriptive and
Justificatory Questions about Induction}
\maketitle

\begin{abstract}
We give a precise account of ``scientific objectivity'' that keeps the
\emph{descriptive} question about inductive reasoning (quid facti?) and the
\emph{normative} question (quid juris?) strictly distinct. The account
identifies objectivity with the property of an inductive procedure's having a
\emph{justification that is invariant under the substitution of any rational
agent for any other} --- i.e.\ with the inter-subjective validity of the
warrant, not with the warrant's actual acceptance. We prove that this account
satisfies the three required adequacy conditions (a)--(c), and we rebut the
deflationary alternative which holds that the descriptive/justificatory
distinction is itself misplaced. The argument is structured as a sequence of
definitions, theses, and lemmas, each established without appeal to the
disputed sociological reduction.
\end{abstract}

\section{Formalisation of the Problem}

We first fix the vocabulary so that no later step can silently exchange one
notion for a more familiar one.

\begin{definition}[Inductive procedure]\label{def:procedure}
An \emph{inductive procedure} is a rule $\Pi$ that maps a body of evidence $E$
(a finite set of propositions describing observations) and a hypothesis $h$ to
an epistemic appraisal $\Pi(h \mid E)$, where the appraisal lies in some fixed
appraisal space $A$ (e.g.\ a qualitative ordering ``confirmed / disconfirmed /
neutral'', or a real-valued degree in $[0,1]$). $\Pi$ is purely syntactic\slash
formal in the sense that $\Pi(h\mid E)$ depends only on the logical and
descriptive content of $E$ and $h$, not on who applies $\Pi$.
\end{definition}

\begin{definition}[The descriptive question, quid facti?]\label{def:facti}
The \emph{descriptive question} is: for which procedures $\Pi$ is it the case
that human inquirers, as a matter of empirical fact, actually employ $\Pi$ (or
something extensionally close to $\Pi$) and do so with apparent practical
success? Let
\[
\mathcal{F} \;=\; \{\,\Pi : \Pi \text{ is, as a matter of fact, employed by
successful inquirers}\,\}.
\]
Membership in $\mathcal{F}$ is a contingent, empirical, psychological\slash
sociological matter, settled by observation of behaviour.
\end{definition}

\begin{definition}[The normative question, quid juris?]\label{def:juris}
The \emph{normative question} is: for which procedures $\Pi$ is there a valid
\emph{justification} $J$ --- a chain of reasons --- showing that adopting $\Pi$
is rationally warranted relative to the goal of forming true (or
truth-conducive) beliefs? Let
\[
\mathcal{J} \;=\; \{\,\Pi : \text{there exists a valid justification } J
\text{ for } \Pi\,\}.
\]
Membership in $\mathcal{J}$ is a logical\slash epistemological matter, settled
by the assessment of reasons, not by observation of behaviour.
\end{definition}

The two questions are distinct because $\mathcal{F}$ and $\mathcal{J}$ are
defined by entirely different membership criteria (behaviour versus warrant).
We do \emph{not} assume in advance that $\mathcal{F} = \mathcal{J}$,
$\mathcal{F} \subseteq \mathcal{J}$, or $\mathcal{J} \subseteq \mathcal{F}$;
indeed the whole point is to keep these inclusions open and answerable
separately.

\begin{definition}[Justification; warrant]\label{def:just}
A \emph{justification} for $\Pi$ is a finite argument whose premises are either
(i) logical or analytic truths, (ii) explicitly adopted methodological
postulates (e.g.\ a uniformity assumption, a simplicity ordering, an admissible
prior), or (iii) further premises themselves justified in the same sense, and
whose conclusion is that $\Pi$ is warranted relative to a stated cognitive aim.
A justification is \emph{valid} iff its inference steps are truth-preserving
(deductively) or warrant-transferring (where the relevant logic of partial
support is itself antecedently fixed).
\end{definition}

\begin{definition}[Rational agent; admissibility]\label{def:agent}
A \emph{rational agent} is any subject who (i) reasons in accordance with the
fixed logic of full and partial support presupposed in
Definition~\ref{def:just}, and (ii) is bound by no premises other than those of
type (i)--(iii) in Definition~\ref{def:just}. A premise or inference is
\emph{admissible} iff it is of one of these types; in particular, ``it is in
fact accepted by community $C$'' is \emph{not} admissible as a premise of a
justification, though it is a perfectly good datum for the descriptive
question.
\end{definition}

\section{The Proposed Account of Objectivity}

\begin{definition}[Inter-subjective invariance of warrant]\label{def:invariance}
A justification $J$ for $\Pi$ is \emph{inter-subjectively invariant} iff its
admissible premises and inference steps contain no ineliminable indexical or
agent-relative element: replacing the agent who advances $J$ by any other
rational agent (Definition~\ref{def:agent}) leaves the validity of $J$
unchanged. Equivalently, $J$ adverts only to the content of $E$, $h$, the
adopted methodological postulates, and the fixed logic of support --- never to
the identity, location, history, or community-membership of the reasoner.
\end{definition}

\begin{definition}[Scientific objectivity --- the proposed account]
\label{def:objectivity}
A procedure $\Pi$ is \emph{objective} iff
\[
\Pi \in \mathcal{O} \;:\Longleftrightarrow\;
\text{there exists an inter-subjectively invariant valid justification } J
\text{ for } \Pi .
\]
A scientific result obtained by applying $\Pi$ to evidence $E$ is
\emph{objective} iff $\Pi \in \mathcal{O}$ and the application of $\Pi$ to $E$
respects the conditions under which $J$ certifies $\Pi$.
\end{definition}

The single, central conceptual move is this: \emph{objectivity is a modal\slash
logical property of the warrant} (does an invariant justification exist?),
\emph{not a statistical property of the population of inquirers} (do people
accept it?). Objectivity is a sub-question of the quid juris? side, sharpened
by the invariance requirement of Definition~\ref{def:invariance}.

\section{Adequacy: Conditions (a)--(c)}

We now prove that Definition~\ref{def:objectivity} meets the three demands of
the problem. Throughout, ``accepted practice'' denotes membership in
$\mathcal{F}$ (Definition~\ref{def:facti}).

\begin{proposition}[Condition (a): no reduction to description]
\label{prop:a}
$\mathcal{O}$ is not definable from, and is not extensionally identical to,
$\mathcal{F}$.
\end{proposition}

\begin{proof}
We exhibit, in principle, members of each of the four cells of the partition
generated by the two predicates ``$\Pi \in \mathcal{O}$'' and ``$\Pi \in
\mathcal{F}$,'' which suffices to show logical independence.

\emph{(i) $\Pi \in \mathcal{F}$ but $\Pi \notin \mathcal{O}$.} A procedure may
be widely and successfully \emph{used} while lacking any inter-subjectively
invariant justification. Example schema: a procedure that privileges
predicates definable only by reference to a particular community's
classificatory habits (a ``parochial'' inductive policy in the style of
Goodman's grue predicates relativised to a tribe). Such a policy can be in
$\mathcal{F}$ --- the tribe uses it and gets by --- yet any candidate
justification must cite the tribe's habits, hence is not invariant
(Definition~\ref{def:invariance}); so $\Pi \notin \mathcal{O}$.

\emph{(ii) $\Pi \in \mathcal{O}$ but $\Pi \notin \mathcal{F}$.} A procedure may
possess an invariant justification yet never be adopted, e.g.\ a correct but
computationally onerous Bayesian conditionalisation under an explicitly
motivated invariant prior, ignored by all actual inquirers for reasons of
cognitive economy. It is in $\mathcal{O}$ by the existence of $J$, but absent
from $\mathcal{F}$.

\emph{(iii) and (iv)} Overlap ($\Pi \in \mathcal{O}\cap\mathcal{F}$) and joint
absence are clearly possible.

Since the four cells are all non-empty in principle, $\mathcal{O} \neq
\mathcal{F}$ and neither inclusion holds by definition. Moreover the
\emph{membership criterion} for $\mathcal{O}$ (existence of an invariant valid
justification) makes no reference to behavioural facts at all
(Definitions~\ref{def:just}--\ref{def:objectivity}); hence $\mathcal{O}$ is not
\emph{definable} from $\mathcal{F}$. This establishes (a): objectivity does not
reduce to a sociological or psychological description of accepted practice.
\end{proof}

\begin{proposition}[Condition (b): the justificatory question is preserved]
\label{prop:b}
The account presupposes, rather than dismisses, the legitimacy of the quid
juris? question.
\end{proposition}

\begin{proof}
By Definition~\ref{def:objectivity}, $\Pi \in \mathcal{O}$ is defined
\emph{via} the existence of a valid justification $J$
(Definition~\ref{def:just}). Hence the notion of objectivity is literally
unintelligible unless the justificatory question ``is $\Pi \in \mathcal{J}$?''
is a genuine, answerable question. Indeed $\mathcal{O} \subseteq \mathcal{J}$:
every objective procedure is a fortiori justifiable, since an
inter-subjectively invariant valid justification is in particular a valid
justification. Thus far from declaring the justificatory question misplaced,
the account makes a refined version of that very question
($\,$existence of an \emph{invariant} warrant$\,$) constitutive of
objectivity. This establishes (b).
\end{proof}

\begin{proposition}[Condition (c): no conflation of classification and
justification]\label{prop:c}
The account secures objectivity without identifying the factual classification
of procedures (membership in $\mathcal{F}$) with the logical problem of their
justification (membership in $\mathcal{J}$, or in $\mathcal{O}$).
\end{proposition}

\begin{proof}
The classificatory task --- sorting which $\Pi$ are in $\mathcal{F}$ --- is
discharged entirely by empirical inquiry into behaviour
(Definition~\ref{def:facti}) and never enters the membership condition for
$\mathcal{O}$. The justificatory task --- sorting which $\Pi$ are in
$\mathcal{J}$, and which of those are in $\mathcal{O}$ --- is discharged
entirely by the assessment of reasons (Definitions~\ref{def:juris},
\ref{def:just}, \ref{def:objectivity}) and never appeals to behaviour. The two
tasks use disjoint evidential bases (observation of practice vs.\ evaluation of
arguments) and disjoint membership criteria. By
Proposition~\ref{prop:a} the resulting sets are logically independent, so no
inference of the form ``$\Pi\in\mathcal{F}$, therefore $\Pi\in\mathcal{O}$''
(or its converse) is licensed. Objectivity is nonetheless secured, because
$\mathcal{O}$ is non-empty whenever some procedure carries an invariant
warrant --- a condition that can obtain (and be recognised to obtain) without
counting heads. This establishes (c).
\end{proof}

\begin{thesis}[Adequacy]\label{thesis:adequacy}
Definition~\ref{def:objectivity} provides a notion of scientific objectivity
satisfying (a), (b), and (c).
\end{thesis}

\begin{proof}
Immediate from Propositions~\ref{prop:a}, \ref{prop:b}, \ref{prop:c}.
\end{proof}

\section{Defence against the Deflationary Alternative}

We now confront the opposing view.

\begin{objection}[Deflationism about the distinction]\label{obj:deflation}
The descriptive/justificatory distinction is itself misplaced. There is no
``view from nowhere''; to call a procedure justified is just to report that it
is endorsed within our best current practice. Hence quid facti? and quid
juris? collapse: the only available answer to ``is $\Pi$ warranted?'' is
``$\Pi$ is what we, the successful inquirers, in fact do.'' Consequently
$\mathcal{O}$, $\mathcal{J}$, and $\mathcal{F}$ coincide, and
Definition~\ref{def:objectivity} mislabels a sociological fact as a logical
property.
\end{objection}

We refute Objection~\ref{obj:deflation} by three independent arguments. The
first two are dilemmas internal to the deflationist's own claim; the third
shows the distinction is indispensable on pain of losing critical purchase.

\begin{lemma}[Self-application dilemma]\label{lem:selfapp}
The deflationist thesis $D$ --- ``warrant $=$ endorsement by successful
practice'' --- cannot be advanced without either (i) covertly invoking the very
distinction it denies, or (ii) self-undermining.
\end{lemma}

\begin{proof}
Ask of $D$ itself: is $D$ warranted? Two cases.

\emph{Case 1: $D$ is offered as warranted in the deflationist's own sense}, i.e.\
``$D$ is endorsed by successful practice.'' Then $D$'s correctness is a mere
contingent report about a community. It then provides no reason for anyone
outside that community --- in particular no reason for the non-deflationist ---
to accept $D$; it is a description of the deflationist's habits, not an argument.
By the deflationist's own lights there is nothing further to say, so $D$ has no
normative force against the proposed account. The objection thereby
\emph{deflates itself}: it reduces to ``we deflationists are disposed to talk
this way,'' which is a datum for $\mathcal{F}$, not a refutation in
$\mathcal{J}$.

\emph{Case 2: $D$ is offered as warranted simpliciter}, i.e.\ as something the
interlocutor \emph{ought} to accept on the basis of reasons, independently of
who happens to endorse it. Then the deflationist is, in the very act of
asserting $D$, presupposing an agent-invariant notion of warrant
(Definition~\ref{def:invariance}) --- exactly the notion whose intelligibility
$D$ denies. The thesis is pragmatically self-refuting.

Either case is fatal: in Case~1 $D$ has no force; in Case~2 $D$ refutes itself.
\end{proof}

\begin{lemma}[Non-coincidence under any non-trivial success criterion]
\label{lem:noncoincide}
If ``success'' in Definition~\ref{def:facti} is given any content beyond ``is
endorsed,'' then $\mathcal{F} \neq \mathcal{J}$ is at least possible, so the
identification $\mathcal{F}=\mathcal{J}$ asserted by the deflationist is not
analytic and stands in need of the very argument it claims to render
unnecessary.
\end{lemma}

\begin{proof}
The deflationist needs ``successful practice,'' not merely ``practice,'' lest
$D$ endorse every fad. But ``success'' must then be cashed out by some
standard --- predictive accuracy, instrumental reliability, coherence, etc.
Call this standard $S$. Now $S$ is either (i) itself reducible to ``what the
community endorses,'' in which case the success qualifier is vacuous and we are
back to bare endorsement (handled by Lemma~\ref{lem:selfapp}, Case~1), or (ii)
$S$ has content independent of endorsement. In case (ii), whether a given
endorsed $\Pi$ \emph{actually meets} $S$ is a further question not settled by
the fact of endorsement: a community can endorse a procedure that fails $S$
(self-deception, lucky streaks, biased sampling of its own track record) and
can fail to endorse a procedure that meets $S$. Hence
$\mathcal{F}\neq\mathcal{J}$ is possible, so their identity is a substantive
thesis, contradicting the claim that the distinction is ``misplaced'' (i.e.\
not even a real question). Thus the deflationist must, after all, argue across
the distinction.
\end{proof}

\begin{lemma}[Indispensability for criticism]\label{lem:criticism}
Collapsing $\mathcal{J}$ into $\mathcal{F}$ makes the proposition ``accepted
scientific practice is mistaken'' contradictory, contrary to a fixed point of
scientific self-understanding; hence the distinction must be retained.
\end{lemma}

\begin{proof}
Under $D$, ``$\Pi$ is unwarranted'' means ``$\Pi$ is not endorsed by successful
practice.'' Then for any actually endorsed-and-successful $\Pi$, ``$\Pi$ is
unwarranted'' is, by $D$, false by definition, and ``current practice is in
error'' is incoherent. But the history of science is in large part the history
of successful, endorsed procedures later judged unwarranted on reflective
grounds (e.g.\ ad hoc rescue strategies, suppressed-evidence inferences,
question-begging uses of auxiliary hypotheses). The intelligibility of such
\emph{immanent self-correction} requires a standard of warrant not constituted
by current endorsement --- precisely $\mathcal{J}\setminus(\text{endorsement})$.
The proposed account supplies this standard via
Definition~\ref{def:invariance}: a community can ask of its own endorsed $\Pi$
whether $\Pi$ carries an agent-invariant warrant, and can discover that it does
not. Deflationism forecloses this question and so cannot represent ordinary
scientific criticism; the distinction is therefore indispensable.
\end{proof}

\begin{thesis}[Refutation of the alternative]\label{thesis:refute}
The view that the descriptive/justificatory distinction is misplaced is
untenable, and the account of Definition~\ref{def:objectivity} is to be
preferred.
\end{thesis}

\begin{proof}
By Lemma~\ref{lem:selfapp} the deflationist thesis either lacks force or
refutes itself. By Lemma~\ref{lem:noncoincide} the identification
$\mathcal{F}=\mathcal{J}$ on which the objection rests is not analytic but
substantive, so the distinction is a genuine question after all. By
Lemma~\ref{lem:criticism} retaining the distinction is required to make sense of
scientific self-criticism. Hence the objection fails on all three counts. The
proposed account, by Thesis~\ref{thesis:adequacy}, meets (a)--(c) while keeping
the two questions separate; it is therefore adequate and preferable to the
deflationary alternative.
\end{proof}

\section{Summary}

Scientific objectivity is the existence, for an inductive procedure, of a
\emph{valid justification invariant under exchange of any rational agent for
any other} (Definitions~\ref{def:invariance}--\ref{def:objectivity}). This:
(a) does not reduce to description, since $\mathcal{O}$ is logically
independent of the behavioural set $\mathcal{F}$
(Proposition~\ref{prop:a}); (b) does not dismiss the justificatory question,
since $\mathcal{O}\subseteq\mathcal{J}$ and the existence of a warrant is
constitutive of objectivity (Proposition~\ref{prop:b}); and (c) does not
conflate classification with justification, since the two are settled by
disjoint evidential bases and criteria (Proposition~\ref{prop:c}). The
deflationary denial of the distinction is refuted by a self-application dilemma,
by the non-analyticity of $\mathcal{F}=\mathcal{J}$, and by the indispensability
of an endorsement-independent standard for scientific self-correction
(Lemmas~\ref{lem:selfapp}--\ref{lem:criticism}).
\end{document}
